\documentclass[11pt]{article}

\usepackage[T1]{fontenc}
\usepackage[utf8]{inputenc}
\usepackage[english]{babel}
\usepackage[margin=1in]{geometry}
\usepackage{microtype}
\usepackage[default,scale=0.95]{sourcesanspro}
\renewcommand{\familydefault}{\sfdefault}

\title{\textbf{Agentic AI for Economic Research:\\ Concepts, Tools, Workflows, and Case Studies}}
\author{Marco Le Moglie\\Catholic University of the Sacred Heart}
\date{1 July 2026}

\begin{document}
\maketitle

\noindent
This four-hour minicourse introduces doctoral students and faculty members to
the practical and critical use of agentic artificial intelligence in economics
research. The course is designed as a research workflow map rather than a
technical manual: its goal is to help participants understand what can be
delegated to AI agents, what must remain under human control, and how to design
AI-assisted research processes that remain transparent, reproducible, and
methodologically sound.

\medskip
\noindent
The first part develops the core vocabulary and conceptual framework. It
explains the shift from chatbots that answer prompts to agents that can inspect
files, call tools, use apps, work inside IDEs or command-line environments, and
iterate through plan--act--verify loops. Participants will learn the practical
meaning of context, tokens, tool use, permissions, sandboxing, autonomy levels,
rules, persistent memory, sub-agents, MCP, CLI and API workflows. Special
attention is given to the changing boundary between chat apps, desktop agent
apps, IDE agents, browser/research assistants, and local or command-line
agents. The course also discusses knowledge bases and memory systems, including
why a curated Obsidian-style Markdown knowledge base can be preferable to a
pure RAG setup for research tasks where provenance, auditability, and human
curation matter.

\medskip
\noindent
The first part also covers the governance side of agentic work: pricing and
token consumption, data protection, institutional constraints, reproducibility,
and the risks associated with hallucinations, automation bias, AI sloppiness,
weak validation, and excessive trust in polished outputs. A central theme is
that agentic AI is safest and most useful when verification is designed into
the workflow through plans, diffs, logs, tests, source checks, and explicit
human review gates.

\medskip
\noindent
The second part turns to live case studies from current economics research
projects. Through practical examples, participants will examine how agents can
support literature review and source triage, data construction and codebook
documentation, Stata/R/Python debugging, replication workflows, robustness
checks, table creation from data, paper rewriting, slide production, grant and
referee-response preparation, and teaching material design. These examples are
used not as generic demonstrations, but as opportunities to understand what
different project difficulties teach about setup, prompting, tool choice,
memory, verification, and when to stop delegating.

\medskip
\noindent
By the end of the minicourse, participants will be able to evaluate agentic AI
tools critically, choose appropriate interfaces for different research tasks,
recognize the privacy and cost implications of agentic workflows, and design
AI-assisted research processes in which human expertise remains responsible for
theory, identification, measurement, interpretation, and final scientific
judgement.

\end{document}
